\documentclass{article}
\usepackage{amsmath,amssymb,amsthm}
\IfFileExists{hyperref.sty}{%
  \IfFileExists{color.sty}{\usepackage{color}\definecolor{paperlinkblue}{rgb}{0.10,0.25,0.42}%
    \usepackage[colorlinks=true,linkcolor=paperlinkblue]{hyperref}}{\usepackage[hidelinks]{hyperref}}%
}{}
\ifdefined\hyperlink
  \newcommand{\paperanchor}[1]{\hypertarget{#1}{}}%
  \newcommand{\paperlink}[2]{\hyperlink{#1}{\begin{NoHyper}#2\end{NoHyper}}}%
\else
  \ifdefined\pdfextension
    \ifnum\outputmode>0
      \def\paperdest{\pdfextension dest}\def\paperstart{\pdfextension startlink}%
      \def\paperend{\pdfextension endlink}\def\papercolor{\pdfextension colorstack}%
    \fi
  \else\ifdefined\pdfstartlink
    \ifnum\pdfoutput>0
      \let\paperdest\pdfdest\let\paperstart\pdfstartlink\let\paperend\pdfendlink\let\papercolor\pdfcolorstack
    \fi
  \fi\fi
  \ifdefined\paperdest
    \newcommand{\paperanchor}[1]{\paperdest name{#1} xyz}%
    \newcommand{\paperlink}[2]{\leavevmode\paperstart attr{/Border [0 0 0]} goto name{#1}%
      \papercolor 0 push {0.10 0.25 0.42 rg 0.10 0.25 0.42 RG}#2\papercolor 0 pop\paperend}%
  \else\ifdefined\XeTeXversion
    \newcommand{\paperanchor}[1]{\special{pdf:dest (#1) [@thispage /XYZ null @ypos null]}}%
    \newcommand{\paperlink}[2]{\leavevmode\special{pdf:bann << /Subtype /Link /Border [0 0 0] /A << /S /GoTo /D (#1) >> >>}%
      \special{color push rgb 0.10 0.25 0.42}#2\special{color pop}\special{pdf:eann}}%
  \else
    \PackageError{lean-to-latex}{No PDF hyperlink backend available}{Use XeLaTeX, pdfLaTeX, LuaLaTeX, or install hyperref.}%
  \fi\fi
\fi
\newcommand{\paperlineanchor}[1]{\raisebox{10pt}[0pt][0pt]{\paperanchor{#1}}}
\DeclareRobustCommand{\paperresultname}[1]{\paperlineanchor{theorem:\arabic{theorem}}#1}
\newtheorem{theorem}{\paperresultname{Theorem}}
\newcounter{paperpremise}
\begin{document}
\section*{Divisibility and parity in the natural numbers}

Divisibility begins with a factor witness. The same idea explains why divisibility passes through sums and products, and why adding an even increment preserves parity.

\subsection*{Glossary}
\begin{description}
\item[{\normalfont\bfseries Factor witness}] A natural number that exhibits the required product for a divisibility claim. For an expression of evenness as a sum of two equal terms, it is the repeated summand.
\item[{\normalfont\bfseries Even}] Divisible by two. The second section gives the equivalent description as a sum of two equal terms.
\item[{\normalfont\bfseries Inductive hypothesis}] The claim at the current index, assumed only while proving the claim at its successor.
\end{description}

\subsection*{Prerequisites}
\begin{description}
\item[\paperlineanchor{premise:4e61742e6d756c5f6f6e65}\paperlink{premise:4e61742e6d756c5f6f6e65}{P1.}] \refstepcounter{paperpremise}\label{premise:4e61742e6d756c5f6f6e65}\textbf{Multiplication by one on the right.} Let \(n\) be a natural number. Then \(n \cdot 1 = n\).
\item[\paperlineanchor{premise:4e61742e6f6e655f6d756c}\paperlink{premise:4e61742e6f6e655f6d756c}{P2.}] \refstepcounter{paperpremise}\label{premise:4e61742e6f6e655f6d756c}\textbf{Multiplication by one on the left.} Let \(n\) be a natural number. Then \(1 \cdot n = n\).
\item[\paperlineanchor{premise:4e61742e6d756c5f636f6d6d}\paperlink{premise:4e61742e6d756c5f636f6d6d}{P3.}] \refstepcounter{paperpremise}\label{premise:4e61742e6d756c5f636f6d6d}\textbf{Multiplication commutativity.} Let \(n\) and \(m\) be natural numbers. Then \(n \cdot m = m \cdot n\).
\item[\paperlineanchor{premise:4e61742e6d756c5f6173736f63}\paperlink{premise:4e61742e6d756c5f6173736f63}{P4.}] \refstepcounter{paperpremise}\label{premise:4e61742e6d756c5f6173736f63}\textbf{Multiplication associativity.} Let \(n\), \(m\) and \(k\) be natural numbers. Then \(n \cdot m \cdot k = n \cdot (m \cdot k)\).
\item[\paperlineanchor{premise:4e61742e6c6566745f64697374726962}\paperlink{premise:4e61742e6c6566745f64697374726962}{P5.}] \refstepcounter{paperpremise}\label{premise:4e61742e6c6566745f64697374726962}\textbf{Multiplication distributes over addition.} Let \(n\), \(m\) and \(k\) be natural numbers. Then \(n \cdot (m + k) = n \cdot m + n \cdot k\).
\item[\paperlineanchor{premise:4e61742e6164645f636f6d6d}\paperlink{premise:4e61742e6164645f636f6d6d}{P6.}] \refstepcounter{paperpremise}\label{premise:4e61742e6164645f636f6d6d}\textbf{Addition commutativity.} Let \(n\) and \(m\) be natural numbers. Then \(n + m = m + n\).
\item[\paperlineanchor{premise:4e61742e6476645f6164645f6966665f7269676874}\paperlink{premise:4e61742e6476645f6164645f6966665f7269676874}{P7.}] \refstepcounter{paperpremise}\label{premise:4e61742e6476645f6164645f6966665f7269676874}\textbf{Divisibility after adding a known multiple.} Let \(k\), \(m\) and \(n\) be natural numbers. If \(k \mid m\), then \(k \mid n \iff k \mid m + n\).
\item[\paperlineanchor{premise:4e61742e74776f5f6d756c}\paperlink{premise:4e61742e74776f5f6d756c}{P8.}] \refstepcounter{paperpremise}\label{premise:4e61742e74776f5f6d756c}\textbf{Doubling by addition.} Let \(n\) be a natural number. Then \(2 \cdot n = n + n\).
\item[\paperlineanchor{premise:4e61742e6d756c5f7a65726f}\paperlink{premise:4e61742e6d756c5f7a65726f}{P9.}] \refstepcounter{paperpremise}\label{premise:4e61742e6d756c5f7a65726f}\textbf{Multiplication by zero.} Let \(n\) be a natural number. Then \(n \cdot 0 = 0\).
\item[\paperlineanchor{premise:4e61742e6164645f7a65726f}\paperlink{premise:4e61742e6164645f7a65726f}{P10.}] \refstepcounter{paperpremise}\label{premise:4e61742e6164645f7a65726f}\textbf{Adding zero on the right.} Let \(n\) be a natural number. Then \(n + 0 = n\).
\item[\paperlineanchor{premise:4e61742e6d756c5f73756363}\paperlink{premise:4e61742e6d756c5f73756363}{P11.}] \refstepcounter{paperpremise}\label{premise:4e61742e6d756c5f73756363}\textbf{Multiplication by a successor.} Let \(n\) and \(m\) be natural numbers. Then \(n \cdot (m + 1) = n \cdot m + n\).
\item[\paperlineanchor{premise:4e61742e6164645f6173736f63}\paperlink{premise:4e61742e6164645f6173736f63}{P12.}] \refstepcounter{paperpremise}\label{premise:4e61742e6164645f6173736f63}\textbf{Addition associativity.} Let \(n\), \(m\) and \(k\) be natural numbers. Then \(n + m + k = n + (m + k)\).
\end{description}

\section{Divisibility and its closure rules}
The first five results exhibit witnesses directly. The remaining results show how to combine witnesses or use the resulting closure rules.

\begin{theorem}\label{theorem:1}
Let \(a\) be a natural number. Then \(a \mid a\).
\end{theorem}

\begin{proof}
Let \(a\) be a natural number.
By \paperlink{premise:4e61742e6d756c5f6f6e65}{Prerequisite~P\ref{premise:4e61742e6d756c5f6f6e65}}, we have \(a \cdot 1 = a\)~(1).
By applying (1), we obtain \(a = a \cdot 1\)~(2).
With factor witness \(1\), (2) shows that \(a \mid a\).
\end{proof}

\begin{theorem}\label{theorem:2}
Let \(a\) be a natural number. Then \(a \mid 0\).
\end{theorem}

\begin{proof}
Let \(a\) be a natural number.
By computation, we have \(0 = a \cdot 0\)~(1).
With factor witness \(0\), (1) shows that \(a \mid 0\).
\end{proof}

\begin{theorem}\label{theorem:3}
Let \(n\) be a natural number. Then \(1 \mid n\).
\end{theorem}

\begin{proof}
Let \(n\) be a natural number.
By \paperlink{premise:4e61742e6f6e655f6d756c}{Prerequisite~P\ref{premise:4e61742e6f6e655f6d756c}}, we have \(1 \cdot n = n\)~(1).
Reversing the equality in (1) gives \(n = 1 \cdot n\)~(2).
With factor witness \(n\), (2) shows that \(1 \mid n\).
\end{proof}

\begin{theorem}\label{theorem:4}
Let \(a\) and \(b\) be natural numbers. Then \(a \mid a \cdot b\).
\end{theorem}

\begin{proof}
Let \(a\) and \(b\) be natural numbers.
By computation, we have
\begin{equation}\tag{1}
a \cdot b = a \cdot b.
\end{equation}
With factor witness \(b\), (1) shows that \(a \mid a \cdot b\).
\end{proof}

\begin{theorem}\label{theorem:5}
Let \(a\) and \(b\) be natural numbers. Then \(a \mid b \cdot a\).
\end{theorem}

\begin{proof}
Let \(a\) and \(b\) be natural numbers.
By \paperlink{premise:4e61742e6d756c5f636f6d6d}{Prerequisite~P\ref{premise:4e61742e6d756c5f636f6d6d}}, we have \(b \cdot a = a \cdot b\)~(1).
With factor witness \(b\), (1) shows that \(a \mid b \cdot a\).
\end{proof}

\begin{theorem}\label{theorem:6}
Let \(a\), \(b\) and \(c\) be natural numbers. If \(a \mid b\) and \(b \mid c\), then \(a \mid c\).
\end{theorem}

\begin{proof}
Let \(a\), \(b\) and \(c\) be natural numbers.
Assume \(a \mid b\)~(1).
Assume \(b \mid c\)~(2).
By (1), choose a natural number \(d\) such that \(b = a \cdot d\)~(3).
By (2), choose a natural number \(e\) such that \(c = b \cdot e\)~(4).
By \paperlink{premise:4e61742e6d756c5f6173736f63}{Prerequisite~P\ref{premise:4e61742e6d756c5f6173736f63}}, we have \(a \cdot d \cdot e = a \cdot (d \cdot e)\)~(5).
By substituting from (4) and (3) and applying (5), we obtain
\begin{equation}\tag{6}
c = a \cdot (d \cdot e).
\end{equation}
With factor witness \(d \cdot e\), (6) shows that \(a \mid c\).
\end{proof}

\begin{theorem}\label{theorem:7}
Let \(a\), \(b\) and \(c\) be natural numbers. If \(a \mid b\) and \(a \mid c\), then \(a \mid b + c\).
\end{theorem}

\begin{proof}
Let \(a\), \(b\) and \(c\) be natural numbers.
Assume \(a \mid b\)~(1).
Assume \(a \mid c\)~(2).
By (1), choose a natural number \(d\) such that \(b = a \cdot d\)~(3).
By (2), choose a natural number \(e\) such that \(c = a \cdot e\)~(4).
By \paperlink{premise:4e61742e6c6566745f64697374726962}{Prerequisite~P\ref{premise:4e61742e6c6566745f64697374726962}}, we have \(a \cdot (d + e) = a \cdot d + a \cdot e\)~(5).
By substituting from (3) and (4) and applying (5), we obtain
\begin{equation}\tag{6}
b + c = a \cdot (d + e).
\end{equation}
With factor witness \(d + e\), (6) shows that \(a \mid b + c\).
\end{proof}

\begin{theorem}\label{theorem:8}
Let \(a\), \(b\) and \(c\) be natural numbers. If \(a \mid b\), then \(a \mid b \cdot c\).
\end{theorem}

\begin{proof}
Let \(a\) and \(b\) be natural numbers.
Assume \(a \mid b\)~(1).
Let \(c\) be a natural number.
Record the instance of \paperlink{theorem:4}{Theorem~\ref{theorem:4}} with its parameters \((a, b)\) set to \((b, c)\) as (2).
The claim follows from \paperlink{theorem:6}{Theorem~\ref{theorem:6}} with its parameters \((a, b, c)\) set to \((a, b, b \cdot c)\), using (1) and (2).
\end{proof}

\begin{theorem}\label{theorem:9}
Let \(a\), \(b\) and \(c\) be natural numbers. If \(a \mid b\), then \(a \mid c \cdot b\).
\end{theorem}

\begin{proof}
Let \(a\) and \(b\) be natural numbers.
Assume \(a \mid b\)~(1).
Let \(c\) be a natural number.
Record the instance of \paperlink{theorem:5}{Theorem~\ref{theorem:5}} with its parameters \((a, b)\) set to \((b, c)\) as (2).
The claim follows from \paperlink{theorem:6}{Theorem~\ref{theorem:6}} with its parameters \((a, b, c)\) set to \((a, b, c \cdot b)\), using (1) and (2).
\end{proof}

\begin{theorem}\label{theorem:10}
Let \(k\), \(m\) and \(n\) be natural numbers. If \(k \mid n\), then \(k \mid m \iff k \mid m + n\).
\end{theorem}

\begin{proof}
Let \(k\), \(m\) and \(n\) be natural numbers.
Assume \(k \mid n\)~(1).
By \paperlink{premise:4e61742e6164645f636f6d6d}{Prerequisite~P\ref{premise:4e61742e6164645f636f6d6d}}, we have \(m + n = n + m\)~(2).
By \paperlink{premise:4e61742e6476645f6164645f6966665f7269676874}{Prerequisite~P\ref{premise:4e61742e6476645f6164645f6966665f7269676874}}, using (1), we have \(k \mid m \iff k \mid n + m\)~(3).
Rewriting (3) using (2) gives
\[
k \mid m \iff k \mid m + n.
\]
\end{proof}

\section{Even numbers}
Specializing divisibility to the divisor two gives both a witness description of even numbers and methods for proving that sums and products are even.

\begin{theorem}\label{theorem:11}
Let \(n\) be a natural number. Then \(2 \mid n + n\).
\end{theorem}

\begin{proof}
Let \(n\) be a natural number.
By \paperlink{premise:4e61742e74776f5f6d756c}{Prerequisite~P\ref{premise:4e61742e74776f5f6d756c}}, we have \(2 \cdot n = n + n\)~(1).
Reversing the equality in (1) gives
\begin{equation}\tag{2}
n + n = 2 \cdot n.
\end{equation}
With factor witness \(n\), (2) shows that \(2 \mid n + n\).
\end{proof}

\begin{theorem}\label{theorem:12}
Let \(n\) be a natural number. If \(2 \mid n\), then \(\exists k \in \mathbb{N}\; (n = k + k)\).
\end{theorem}

\begin{proof}
Let \(n\) be a natural number.
Assume \(2 \mid n\)~(1).
By (1), choose a natural number \(k\) such that \(n = 2 \cdot k\)~(2).
By \paperlink{premise:4e61742e74776f5f6d756c}{Prerequisite~P\ref{premise:4e61742e74776f5f6d756c}}, we have \(2 \cdot k = k + k\)~(3).
Combining (2) and (3) gives \(n = k + k\)~(4).
With factor witness \(k\), (4) shows that
\[
\exists k_{1} \in \mathbb{N}\; (n = k_{1} + k_{1}).
\]
\end{proof}

\begin{theorem}\label{theorem:13}
Let \(n\) be a natural number. If \(\exists k \in \mathbb{N}\; (n = k + k)\), then \(2 \mid n\).
\end{theorem}

\begin{proof}
Let \(n\) be a natural number.
Assume
\begin{equation}\tag{1}
\exists k \in \mathbb{N}\; (n = k + k).
\end{equation}
By (1), choose a natural number \(k\) such that \(n = k + k\)~(2).
Record the instance of \paperlink{theorem:11}{Theorem~\ref{theorem:11}} with its parameter \(n\) set to \(k\) as (3).
Rewriting (3) using (2) gives \(2 \mid n\).
\end{proof}

\begin{theorem}\label{theorem:14}
Let \(n\) be a natural number. Then \(2 \mid n \iff \exists k \in \mathbb{N}\; (n = k + k)\).
\end{theorem}

\begin{proof}
Let \(n\) be a natural number.
We prove both implications.
\begin{description}
\item[Forward implication.] Assume \(2 \mid n\)~(1).
The claim follows from \paperlink{theorem:12}{Theorem~\ref{theorem:12}} with its parameter \(n\) set to \(n\), using (1).
\item[Reverse implication.] Assume
\begin{equation}\tag{2}
\exists k \in \mathbb{N}\; (n = k + k).
\end{equation}
The claim follows from \paperlink{theorem:13}{Theorem~\ref{theorem:13}} with its parameter \(n\) set to \(n\), using (2).
\end{description}
This proves the equivalence.
\end{proof}

\begin{theorem}\label{theorem:15}
Let \(m\) and \(n\) be natural numbers. If \(2 \mid m\) and \(2 \mid n\), then \(2 \mid m + n\).
\end{theorem}

\begin{proof}
Let \(m\) and \(n\) be natural numbers.
Assume \(2 \mid m\)~(1).
Assume \(2 \mid n\)~(2).
The claim follows from \paperlink{theorem:7}{Theorem~\ref{theorem:7}} with its parameters \((a, b, c)\) set to \((2, m, n)\), using (1) and (2).
\end{proof}

\begin{theorem}\label{theorem:16}
Let \(n\) and \(m\) be natural numbers. If \(2 \mid n\), then \(2 \mid n \cdot m\).
\end{theorem}

\begin{proof}
Let \(n\) and \(m\) be natural numbers.
Assume \(2 \mid n\)~(1).
The claim follows from \paperlink{theorem:8}{Theorem~\ref{theorem:8}} with its parameters \((a, b, c)\) set to \((2, n, m)\), using (1).
\end{proof}

\begin{theorem}\label{theorem:17}
Let \(n\) be a natural number. If \(2 \mid n\), then \(\exists k \in \mathbb{N}\; (n \cdot n = k + k)\).
\end{theorem}

\begin{proof}
Let \(n\) be a natural number.
Assume \(2 \mid n\)~(1).
Record the instance of \paperlink{theorem:8}{Theorem~\ref{theorem:8}} with its parameters \((a, b, c)\) set to \((2, n, n)\), using (1) as (2).
The claim follows from \paperlink{theorem:12}{Theorem~\ref{theorem:12}} with its parameter \(n\) set to \(n \cdot n\), using (2).
\end{proof}

\begin{theorem}\label{theorem:18}
Let \(k\), \(k_{1}\) and \(n\) be natural numbers. If \(2 \mid k\) and \(2 \mid k_{1}\), then \(2 \mid k + k_{1} \cdot n\).
\end{theorem}

\begin{proof}
Let \(k\), \(k_{1}\) and \(n\) be natural numbers.
Assume \(2 \mid k\)~(1).
Assume \(2 \mid k_{1}\)~(2).
We proceed by induction on \(n\).
\begin{description}
\item[Base case.] We must show \(2 \mid k + k_{1} \cdot 0\).
By \paperlink{premise:4e61742e6d756c5f7a65726f}{Prerequisite~P\ref{premise:4e61742e6d756c5f7a65726f}}, we have \(k_{1} \cdot 0 = 0\)~(3).
Applying \(k_{2} \mapsto k + k_{2}\) to both sides of (3) gives
\begin{equation}\tag{4}
k + k_{1} \cdot 0 = k + 0.
\end{equation}
By \paperlink{premise:4e61742e6164645f7a65726f}{Prerequisite~P\ref{premise:4e61742e6164645f7a65726f}}, we have \(k + 0 = k\)~(5).
Combining (4) and (5) gives
\begin{equation}\tag{6}
k + k_{1} \cdot 0 = k.
\end{equation}
Rewriting (1) using (6) gives
\[
2 \mid k + k_{1} \cdot 0.
\]
\item[Inductive step.] Let \(i\) be a natural number.
As the inductive hypothesis, assume
\begin{equation}\tag{7}
2 \mid k + k_{1} \cdot i.
\end{equation}
We must show \(2 \mid k + k_{1} \cdot (i + 1)\).
Record the instance of \paperlink{theorem:7}{Theorem~\ref{theorem:7}} with its parameters \((a, b, c)\) set to \((2, k + k_{1} \cdot i, k_{1})\), using (7) and (2) as (8).
By \paperlink{premise:4e61742e6d756c5f73756363}{Prerequisite~P\ref{premise:4e61742e6d756c5f73756363}}, we have \(k_{1} \cdot (i + 1) = k_{1} \cdot i + k_{1}\)~(9).
Applying \(k_{2} \mapsto k + k_{2}\) to both sides of (9) gives
\begin{equation}\tag{10}
k + k_{1} \cdot (i + 1) = k + (k_{1} \cdot i + k_{1}).
\end{equation}
By \paperlink{premise:4e61742e6164645f6173736f63}{Prerequisite~P\ref{premise:4e61742e6164645f6173736f63}}, we have \(k + k_{1} \cdot i + k_{1} = k + (k_{1} \cdot i + k_{1})\)~(11).
Reversing the equality in (11) gives
\begin{equation}\tag{12}
k + (k_{1} \cdot i + k_{1}) = k + k_{1} \cdot i + k_{1}.
\end{equation}
Combining (10) and (12) gives
\begin{equation}\tag{13}
k + k_{1} \cdot (i + 1) = k + k_{1} \cdot i + k_{1}.
\end{equation}
Rewriting (8) using (13) gives
\[
2 \mid k + k_{1} \cdot (i + 1).
\]
\end{description}
This completes the induction.
\end{proof}

\section{Parity in Arithmetic Progressions}
Adding an even multiple preserves parity. We first establish how divisibility behaves when a divisible quantity is added, then apply this fact to even increments.

\begin{theorem}\label{theorem:19}
Let \(q\), \(n\) and \(k\) be natural numbers. If \(q \mid k\), then \(q \mid n \iff q \mid n + k\).
\end{theorem}

\begin{proof}
Let \(q\), \(n\) and \(k\) be natural numbers.
Assume \(q \mid k\)~(1).
The claim follows from \paperlink{theorem:10}{Theorem~\ref{theorem:10}} with its parameters \((k, m, n)\) set to \((q, n, k)\), using (1).
\end{proof}

\begin{theorem}\label{theorem:20}
Let \(q\), \(n\), \(k\) and \(k_{1}\) be natural numbers. If \(q \mid k\), then \(q \mid n \iff q \mid n + k \cdot k_{1}\).
\end{theorem}

\begin{proof}
Let \(q\), \(n\), \(k\) and \(k_{1}\) be natural numbers.
Assume \(q \mid k\)~(1).
Record the instance of \paperlink{theorem:8}{Theorem~\ref{theorem:8}} with its parameters \((a, b, c)\) set to \((q, k, k_{1})\), using (1) as (2).
The claim follows from \paperlink{theorem:10}{Theorem~\ref{theorem:10}} with its parameters \((k, m, n)\) set to \((q, n, k \cdot k_{1})\), using (2).
\end{proof}

\begin{theorem}\label{theorem:21}
Let \(n\), \(d\) and \(k\) be natural numbers. Then \(2 \mid n \iff 2 \mid n + (d + d) \cdot k\).
\end{theorem}

\begin{proof}
Let \(n\), \(d\) and \(k\) be natural numbers.
Record the instance of \paperlink{theorem:11}{Theorem~\ref{theorem:11}} with its parameter \(n\) set to \(d\) as (1).
The claim follows from \paperlink{theorem:20}{Theorem~\ref{theorem:20}} with its parameters \((q, n, k, k_{1})\) set to \((2, n, d + d, k)\), using (1).
\end{proof}

\begin{theorem}\label{theorem:22}
Let \(n\) and \(k\) be natural numbers. If \(2 \mid n\), then \(2 \mid n \cdot k\).
\end{theorem}

\begin{proof}
Let \(n\) and \(k\) be natural numbers.
Assume \(2 \mid n\)~(1).
The claim follows from \paperlink{theorem:8}{Theorem~\ref{theorem:8}} with its parameters \((a, b, c)\) set to \((2, n, k)\), using (1).
\end{proof}

\begin{theorem}[Main result]\label{theorem:23}
Let \(n\), \(k\) and \(k_{1}\) be natural numbers. If \(2 \mid k\), then \(2 \mid n \iff 2 \mid n + k \cdot k_{1}\).
\end{theorem}

\begin{proof}
Let \(n\), \(k\) and \(k_{1}\) be natural numbers.
Assume \(2 \mid k\)~(1).
The claim follows from \paperlink{theorem:20}{Theorem~\ref{theorem:20}} with its parameters \((q, n, k, k_{1})\) set to \((2, n, k, k_{1})\), using (1).
\end{proof}

\end{document}
